\section{Objetivos}

  A partir de las deficiencias identificadas en WebSIS y de las necesidades de los estudiantes de la UMSS, se plantean los siguientes objetivos.

  \subsection{Objetivo general}

    Analizar las barreras de usabilidad, arquitectura de información y diseño de interacción presentes en WebSIS, y proponer una solución interactiva centrada en el estudiante que facilite la gestión académica digital dentro del sistema con mayor claridad, autonomía y confianza.

  \subsection{Objetivos específicos}

    \begin{itemize}
      \item Mapear los módulos estudiantiles de WebSIS, especialmente inscripción, notas, horarios y situación académica.

      \item Identificar y evaluar las necesidades y puntos de fricción de los estudiantes y  el estado actual de la interfaz, mediante encuestas, entrevistas y observación, aplicando heurísticas de usabilidad y clasificando los problemas por severidad.

      \item Analizar y diseñar la nueva arquitectura de información y su correspondencia con el modelo mental del estudiante.

      \item Diseñar prototipos interactivos de baja fidelidad que atienda los problemas críticos identificados.

      \item Justificar el rediseño con base en usabilidad, leyes de la Gestalt y reducción de carga cognitiva.
    \end{itemize}
